%% kapa: exam=1 round=29 year=2018 period=1 form=A subject=civil_law subject_ko=민법 seq=20 number=20
%% source: Q-Net 감정평가사 기출문제 (cst00309) artlSeq=5202808 file=extracted/1008/5202808/2018년 제29회 감정평가사 1차시험 기출문제/2018년 제29회 감정평가사 1차시험 기출문제/1교시 A형.pdf
%% conversion: kordoc 4.24.51 -> tools/gen1.py (auto); review=proofread-sonnet; flags=none
\begin{question}{20}
기간의 계산에 관한 설명으로 옳지 않은 것은? (기간 말일의 공휴일 등 기타 사유는 고려하지 않음)
\begin{choices}
  \item 기간을 연으로 정한 경우 최종의 월에 해당일이 없는 때에는 그 익월의 초일로 기간이 만료한다.
  \item 기간을 일(日)로 정한 때에는 기간말일의 종료로 기간이 만료한다.
  \item 기간을 시, 분, 초로 정한 때에는 즉시로부터 기산한다.
  \item 기간을 월로 정한 경우 그 기간이 오전 영시로부터 시작하는 때에는 기간의 초일을 산입한다.
  \item 연령계산에는 출생일을 산입한다.
\end{choices}
\end{question}
\answer{1}
